\section{Step 25 Recognize when the buffer becomes full}\label{sha-s25}

Advance the full source pointer, reduce remaining input, and increase the used count. If the buffer has fewer than sixty-four bytes, continue at update\_byte. Reaching sixty-four must leave the copy loop and enter compression before another store occurs. Otherwise the next byte would overwrite the schedule area following the buffer.

\subsection{Write and inspect}
Replace only the temporary instructions between \texttt{\# BEGIN upd\_walk} and \texttt{\# END upd\_walk}. Keep the supplied labels outside that region.

\Needspace{10\baselineskip}
\MARSCodeLabel{MIPS assembly}
\begin{minted}[samepage=true]{mars-idea}
    daddiu $s1, $s1, 1
    daddiu $s2, $s2, -1
    addiu $s3, $s3, 1
    li $t0, 64
    bne $s3, $t0, update_byte
    nop
\end{minted}

Use recipe \texttt{s25-1} with the cumulative program through this edit. The exact recipe names the entry, prepared input and checkpoint. Observe the checkpoint before continuing to its temporary or normal return.

\subsection{Expected observation}
For a 64-byte update to an empty context, the buffer is full, remaining input is zero and used is sixty-four at the compression boundary.

\Needspace{9\baselineskip}
\subsection{Check your understanding}
Why not immediately copy byte sixty-five into the same buffer?

\begin{itemize}
  \item \textbf{A} It would exceed the 64-byte buffer.
  \item \textbf{B} It would make the message odd-length.
  \item \textbf{C} It would alter the constant table automatically.
\end{itemize}
